\section{USB глубже: CDC, накопитель и роль Host}\label{les:usb-advanced}
\lessonintro
  {различать классы и роли USB и учитывать общий PHY}
  {уроки 19 и 11}
  {канал команд ping/pong через нативный USB CDC}
  {USB-кабель; для дополнительного опыта --- подготовленная host-плата}

\subsection{Класс и роль --- разные понятия}
Класс USB описывает данные: HID --- ввод, CDC --- последовательный канал,
MSC --- блочное хранилище. Composite-устройство объединяет интерфейсы в одном наборе дескрипторов.
VID/PID и строки описания --- часть интерфейса устройства, а не случайные числа.
В роли Device плата отвечает компьютеру. В роли Host сама перечисляет подключённое устройство
и управляет обменом. Поддержка одного класса не означает поддержку любого USB-устройства.

ESP32-S3 поддерживает Full Speed: 12 Мбит/с на линии, а не обещанную скорость файлов.
Внутренний PHY общий для USB-OTG и фиксированного USB Serial/JTAG. На обычной плате
нельзя получить одновременно независимый TinyUSB и встроенную USB-JTAG-отладку через те же D+/D−.
Мост UART остаётся независимым путём прошивки и диагностики.

\subsection{Своя консоль TinyUSB CDC}
Выберите USB Mode: USB-OTG (TinyUSB), USB CDC On Boot: Disabled.
Объект CDC создаёт сам скетч, поэтому автоматический CDC отключён.
Подключите нативный USB, откройте появившийся порт и отправьте \code{ping} с переводом строки.
При этих настройках диагностический \code{Serial} направлен в мост UART.
Поэтому в этом опыте к компьютеру подключают \emph{оба} разъёма DevKitC-1.

\begin{figure}[H]
\centering
\begin{tikzpicture}[font=\sffamily\scriptsize]
  \node[blk,fill=blkper,minimum width=2.2cm,minimum height=3cm] (pc) at (0,0) {Компьютер};
  \node[blk,fill=blkrtc,minimum width=1.6cm,minimum height=0.7cm] (c1) at (4.2,0.9) {«UART»};
  \node[blk,fill=blkrtc,minimum width=1.6cm,minimum height=0.7cm] (c2) at (4.2,-0.9) {«USB»};
  \node[blk,fill=blkmem,minimum width=2.3cm,minimum height=1.0cm] (br) at (7.6,0.9) {Мост\\ USB--UART};
  \node[blk,fill=blkcpu,minimum width=2.6cm,minimum height=3cm] (esp) at (11.4,0) {};
  \node[font=\sffamily\small\bfseries] at (11.4,1.1) {ESP32-S3};
  \node[align=center] at (11.4,0) {UART0\\ (\pin{43}, \pin{44})};
  \node[align=center] at (11.4,-0.95) {USB: D$-$ \pin{19},\\ D$+$ \pin{20}};
  \draw[thick,accent] (pc.east |- c1) -- node[above]{кабель 1} (c1.west);
  \draw[thick,esp] (pc.east |- c2) -- node[above]{кабель 2} (c2.west);
  \draw[thick,accent] (c1.east) -- (br.west);
  \draw[thick,accent] (br.east) -- (esp.west |- br.east);
  \draw[thick,esp] (c2.east) -- node[above]{D$-$, D$+$} (esp.west |- c2.east);
  \node[align=center,text=accent] at (5.9,1.75) {прошивка, \code{Serial}};
  \node[align=center,text=esp] at (6.4,-1.65) {TinyUSB CDC: команды \code{ping} / \code{pong}};
\end{tikzpicture}
\caption{Какой разъём для чего: «UART» через мост даёт прошивку и диагностический \code{Serial},
«USB» идёт прямо на выводы \pin{19}/\pin{20} и несёт TinyUSB CDC}
\end{figure}

\codecap{Одна команда на строку через USB CDC}
\codeinput{l34_usb_cdc}

CDC передаёт поток байтов: команда может прийти частями, несколько команд --- одним чтением.
Поэтому выделяем строку самостоятельно и ограничиваем её длину. Значение baud rate виртуального
CDC не задаёт физическую скорость USB так же, как у UART. Опишите кодировку, завершение строки,
максимальную длину и ответы на неизвестные команды. Проверьте отключение и повторное подключение:
программа не должна бесконечно ждать открытого монитора.

\subsection{MSC: кто владеет файловой системой}
Накопитель предъявляет хосту сектора, а не готовые файлы. FAT и каталоги обслуживает хост
или файловая библиотека. Если компьютер и ESP одновременно изменяют одну файловую систему,
их кэши могут разойтись и повредить данные. Назначьте одного владельца: пока диск смонтирован
на компьютере, приложение туда не пишет; после извлечения владельцем может стать плата.
Первый дополнительный опыт делайте с отдельным тестовым RAM-диском, не с картой ученических работ.
Пример CDC диски не создаёт и файлы не изменяет.

\subsection{Host: переходника мало}
Хост подаёт устройству 5 В VBUS через рассчитанную цепь с ограничением тока и не конфликтует
с питанием компьютера. USB-переходник сам не добавляет эту схему в DevKitC-1.
Используйте host-плату или документированную схему с управляемым питанием.
Далее нужны host-стек и драйвер класса: ESP-IDF содержит примеры HID, CDC и MSC.
Первый опыт --- чтение дескрипторов собственной клавиатуры и вывод событий в UART.
Снятие устройства обрабатывается как штатное событие, а не бесконечное ожидание данных.

\begin{trouble}
Порт не появился: разъём, режим или кабель. После прошивки порт изменился: обновите выбор.
JTAG пропал при TinyUSB: контроллеры используют общий PHY. Восстанавливайте доступ через
мост и BOOT/RST; eFuse для этих опытов не меняют.
\end{trouble}
\begin{summary}
Класс описывает данные, роль --- инициатора обмена. CDC требует выделения команд,
MSC --- одного владельца диска, Host --- корректного VBUS и драйвера класса.
\end{summary}

\begin{selfcheck}
\begin{enumerate}
\item Может ли одна команда CDC прийти несколькими чтениями?
\item Почему нельзя одновременно писать на диск с двух сторон?
\item Почему переходника недостаточно для Host?
\end{enumerate}
\end{selfcheck}
\needspace{20\baselineskip}
\begin{tasks}
\begin{enumerate}
\item \lv{1} Укажите роли в парах «компьютер — плата» и «плата — клавиатура».
\item \lv{1} Получите pong на команду ping.
\item \lv{2} Добавьте uptime с временем в миллисекундах.
\item \lv{2} Проверьте команды частями и две команды вместе, задайте реакцию на переполнение.
\item \lv{3} Составьте автомат владения MSC-диском: приложение, хост, извлечение, возврат.
\item \lv{4} На host-плате прочитайте дескрипторы своей клавиатуры и обработайте отключение без зависания.
\end{enumerate}
\end{tasks}

\paragraph{Документация для продолжения}\href{https://docs.espressif.com/projects/esp-idf/en/release-v5.5/esp32s3/api-guides/usb-serial-jtag-console.html}{Официальное руководство по теме}.
\href{https://docs.espressif.com/projects/esp-idf/en/v5.4/esp32s3/api-reference/peripherals/usb_host.html}{USB Host и драйверы классов}.
